\chapter{Conclusão}

O presente trabalho teve como objetivo propor e validar um modelo de integração entre as práticas de DevSecOps e Infraestrutura como Código (IaC), visando a automação da segurança e conformidade em ambientes de computação em nuvem. Diante do cenário atual, onde as ameaças cibernéticas e os custos de violações de dados crescem exponencialmente, a transição de um modelo de segurança reativo para uma estratégia de Shift-left mostrou-se não apenas recomendável, mas essencial para a sustentabilidade das operações de TI.

A Prova de Conceito desenvolvida demonstrou que a convergência tecnológica entre Terraform, AWS e ferramentas de análise estática como Checkov, GitGuardian e OPA é perfeitamente viável e altamente eficaz. Os resultados práticos confirmaram que a automação é capaz de eliminar o erro humano em configurações críticas — como a exposição indevida de portas de rede e a ausência de criptografia em dados sensíveis — antes mesmo que o recurso seja provisionado. A detecção de 99\% das falhas críticas durante os testes valida a robustez dos portões de qualidade (quality gates) implementados na esteira de CI/CD.

Um dos principais mitos desmistificados por esta pesquisa é a percepção de que a segurança atua como um gargalo para a agilidade do desenvolvimento. A redução do tempo de revisão de segurança, de uma média manual de 20 minutos para menos de 30 segundos, prova que a automação DevSecOps é um acelerador de entregas seguras. Além disso, o uso do Open Policy Agent (OPA) demonstrou que a esteira de infraestrutura pode ir além da segurança técnica, incorporando regras de governança financeira (FinOps) e conformidade jurídica (LGPD) de forma nativa e transparente.

Apesar da eficácia comprovada, a adoção deste modelo exige uma mudança cultural nas organizações, onde a segurança passa a ser uma responsabilidade compartilhada e o código de infraestrutura é tratado com o mesmo rigor e disciplina de um código de aplicação. A ocorrência de uma baixa taxa de falsos positivos reforça a necessidade de um processo de melhoria contínua e ajuste fino das políticas de segurança para cada contexto organizacional.

Por fim, este estudo cumpre seus objetivos gerais e específicos ao entregar um guia prático para a implementação de infraestruturas resilientes e auditáveis. Como sugestão para trabalhos futuros, recomenda-se a exploração da integração de ferramentas de análise dinâmica (DAST) em ambientes de execução e a aplicação de inteligência artificial para a detecção preditiva de anomalias em códigos IaC. A automação da segurança, conforme demonstrado, não é um destino final, mas um processo contínuo de evolução rumo a operações de TI mais seguras, eficientes e escaláveis.